\section{Local kinetic defects and the uniform-mobility baseline}
\label{sec:local-baseline}
Throughout this section $D(s)=\epsilon>0$ and
$J_\epsilon(c)=J_{k_c}(\epsilon)$, where
\begin{equation}\label{eq:cosine-ensemble}
 \Gamma=\R/(2\pi\mathbb Z),\qquad k_c(s)=(c+\cos s)^2,
 \qquad c\sim\mathrm{Uniform}[-2,2].
\end{equation}
The budget of this uniform design is $M=2\pi\epsilon$. A simple zero
of $c+\cos s$ gives a quadratic zero of the rate. At $c=\pm1$ two
such zeros merge, giving a quartic zero. We first justify the local
source problems and the uniform estimates needed to average through
these mergers. These are estimates of the exact scalar response;
no bounded scalar error is inferred from a leading equivalent.

\subsection{Whole-line source problems and interval localization}
For a nonnegative confining potential $W$, write
\begin{equation}\label{eq:line-response}
 \mathscr J(W)=\sup_{v\in C_c^\infty(\R)}
 \left\{2\int_\R v-\int_\R(|v'|^2+Wv^2)\right\}.
\end{equation}
The notation $\langle1,(-\partial_x^2+W)^{-1}1\rangle$ below always
means this energy-dual quantity, not an $L^2(\R)$ inverse applied to
an $L^2$ constant. On an interval $I$, let $\mathscr J_I^{\rm D}(W)$
and $\mathscr J_I^{\rm N}(W)$ denote the same supremum over
$H_0^1(I)$ and $H^1(I)$, respectively. The superscripts specify the
variational spaces; the Neumann condition is natural in the weak equation.

\begin{lemma}[Expanding intervals, including free endpoints]
\label{lem:interval-response}
For $W(x)=x^2+z$, with $z$ in a compact subset of $[0,\infty)$, or
$W(x)=(x^2-\mu)^2$, with $\mu$ in a compact subset of $\R$, the
whole-line response is finite and positive. Both interval responses
on $(-R,R)$ converge to it as $R\to\infty$, uniformly in the stated
compact parameter sets. Their source solutions and the whole-line
source solution satisfy, with $p=2$ or $4$ respectively,
\begin{align}
 \int (|v'|^2+Wv^2)&\ge c\int(1+|x|^p)v^2,
                      \label{eq:interval-anchor}\\
 \int_{|x|>L}|v|&\le C L^{(1-p)/2}
       \left[\int (|v'|^2+Wv^2)\right]^{1/2},\qquad L\ge L_0.
                      \label{eq:interval-tail}
\end{align}
The integrals on a finite interval are restricted to that interval;
\eqref{eq:interval-anchor} holds for every $H^1$ function there once
$R\ge R_0$. The constants are uniform over each compact parameter set.

The convergence remains true along varying intervals $R_n\to\infty$
and parameters converging in those sets if the derivative, potential,
and source integrals acquire positive weights uniformly bounded above
and below, with the weights converging locally uniformly to one.
\end{lemma}
\begin{proof}
Choose $L_0$ beyond all potential zeros and large enough that
$W(x)\ge c|x|^p$ for $|x|\ge L_0$, uniformly in the parameters.
On the fixed core $[-L_0-1,L_0+1]$, the fundamental theorem of calculus,
followed by averaging an anchor point in $[L_0,L_0+1]$, gives
\[
 \int_{-L_0}^{L_0}v^2\le C\left(\int_{-L_0-1}^{L_0+1}|v'|^2
                                  +\int_{L_0}^{L_0+1}v^2\right).
\]
The anchor interval has a uniformly positive potential. The potential
controls both the exterior $L^2$ mass and its $|x|^p$ weight. This proves
\eqref{eq:interval-anchor}, without endpoint conditions. Weighted
Cauchy--Schwarz proves \eqref{eq:interval-tail}, since
$\int_{|x|>L}|x|^{-p}\dd x=2L^{1-p}/(p-1)$. It also proves
$|\int v|\le C\mathfrak e[v]^{1/2}$ for the indicated energy
$\mathfrak e$. Thus the source is continuous on its Hilbert energy
completion, and Riesz representation gives a unique maximizer $h$ with
$\mathfrak e[h]=\int h=\mathscr J$. It is nonnegative by replacing a
trial with its absolute value. The same argument on finite intervals
gives uniformly bounded responses and energies of their maximizers.
A nonzero nonnegative compact test proves positivity.

Compact smooth functions are dense in the whole-line energy space:
cut off first, using the integrability of $v',v$, and $|x|^{p/2}v$,
and then mollify on the compact support. Every such test is available
on every sufficiently large Dirichlet interval. This proves the lower
limit for both interval problems. For the Neumann upper limit, take a
sequence of maximizers realizing the upper limit. The energy bound and
\eqref{eq:interval-anchor} give a subsequence weakly convergent in
$H^1$ and strongly convergent in $L^2$ on every compact interval.
The limit $h$ has no greater energy, by lower semicontinuity and then
expansion of the compact interval. Equation \eqref{eq:interval-tail}
makes the source tails uniformly small, so the source integrals converge.
Consequently
\[
 \limsup\mathscr J_{(-R,R)}^{\rm N}(W)
 \le2\int h-\mathfrak e[h]\le\mathscr J(W).
\]
The same proof allows converging parameters and locally converging
weights: uniform comparability preserves the anchor and tail bounds,
and local convergence identifies the limiting energy and source.
It proves continuity of the whole-line response by applying the same
argument directly on $\R$. Finally, failure of uniform convergence on
a compact parameter set would give a sequence with a convergent
parameter subsequence contradicting this argument.
\end{proof}

A partition of a circle or interval gives the useful bracketing rule
\begin{equation}\label{eq:response-bracketing}
 \sum_j\mathscr J_{I_j}^{\rm D}\le J\le
 \sum_j\mathscr J_{I_j}^{\rm N}.
\end{equation}
For the lower bound one may retain only selected intervals and extend
their zero-trace trials by zero. For the upper bound independent endpoint
traces enlarge the trial space. Energies and sources are additive, so
there are no interface terms. On a complementary interval where the
potential is positive, discarding the derivative energy bounds its
response by $\int W^{-1}$. The same enlargement bounds the original
lowest Rayleigh quotient below by the minimum of the Neumann quotients
of the pieces.

\begin{proposition}[Harmonic coefficient and isolated zeros]
\label{prop:harmonic}
For $z\ge0$,
\begin{equation}\label{eq:harmonic-constant}
 \mathscr J(x^2+z)=\frac\pi2
     \frac{\Gamma((z+1)/4)}{\Gamma((z+3)/4)},\qquad
 C_0:=\mathscr J(x^2)=\frac{\pi\Gamma(1/4)}{2\Gamma(3/4)}.
\end{equation}
In particular $C_0=4.6474760094\ldots$. If a fixed nonnegative $C^2$
rate $k_0$ on a compact circle has finitely many zeros $s_j$, at least
one, all with
$k_0(s_j+x)=a_jx^2+o(x^2)$ and $a_j>0$, then
\begin{equation}\label{eq:isolated-harmonic}
 J_{k_0}(\epsilon)\sim C_0\epsilon^{-1/4}\sum_j a_j^{-3/4}.
\end{equation}
If there are no zeros, $J_{k_0}(\epsilon)\to\int k_0^{-1}$ instead.
\end{proposition}
\begin{proof}
For completeness the oscillator kernel is
\[
 K_t(x,y)=\frac{1}{\sqrt{2\pi\sinh(2t)}}
 \exp\left[-\frac{(x^2+y^2)\cosh(2t)-2xy}{2\sinh(2t)}\right].
\]
Differentiating verifies $\partial_tK_t=(\partial_x^2-x^2)K_t$;
the small-$t$ Gaussian limit is the identity kernel. Thus it gives
the oscillator heat semigroup. Two Gaussian integrations yield
$\int_{\R^2}K_t=\sqrt{2\pi/\sinh(2t)}$. First apply the $L^2$
resolvent identity to $\ind_{[-L,L]}$. These sources converge to the
constant in the dual energy norm by \eqref{eq:interval-tail}, and
kernel positivity permits monotone convergence. Hence
\[
 \mathscr J(x^2+z)=\int_0^\infty e^{-zt}
             \sqrt{\frac{2\pi}{\sinh(2t)}}\dd t
 =\frac{\sqrt\pi}{2}\Bfun\left(\frac{z+1}{4},\frac12\right),
\]
where the last substitution is $r=e^{-4t}$. This proves
\eqref{eq:harmonic-constant} without treating $1$ as an $L^2$ vector.

Choose disjoint fixed intervals around the zeros with
$(a_j-\eta)x^2\le k_0(s_j+x)\le(a_j+\eta)x^2$.
Their complement has a positive lower bound for $k_0$ and contributes
$O_\eta(1)$ to the upper bracket. Scaling $x=(\epsilon/a)^{1/4}y$
shows that either local quadratic interval response equals
$\epsilon^{-1/4}a^{-3/4}\mathscr J_{(-R,R)}^B(y^2)$, where
$R\to\infty$ and $B\in\{\mathrm D,\mathrm N\}$.
Apply \cref{lem:interval-response}, then let $\eta\downarrow0$.
This proves \eqref{eq:isolated-harmonic}. With no zeros, the upper bound
is $\int k_0^{-1}$; the trial $v=1/k_0$, approximated if necessary in
$H^1$, gives the matching limit as its derivative energy tends to zero.
\end{proof}

\subsection{A merging pair of zeros}
\begin{definition}\label{def:pair-response}
The quartic pair response is
\begin{equation}\label{eq:pair-definition}
 \mathcal C_{\rm pair}(\mu)=\mathscr J((x^2-\mu)^2),\qquad \mu\in\R.
\end{equation}
It is a positive continuous function by \cref{lem:interval-response}.
\end{definition}
For a local rate $(bx^2-t)^2$ with $b>0$ and constant mobility
$\epsilon$, the changes of variables
\begin{equation}\label{eq:fold-scaling}
 \ell=(\epsilon/b^2)^{1/6},\quad
 \mu=\frac{t}{\epsilon^{1/3}b^{1/3}},\quad x=\ell y
\end{equation}
give the exact whole-line response
$\epsilon^{-1/2}b^{-1}\mathcal C_{\rm pair}(\mu)$.
Indeed the energy scale of the operator is
$\epsilon^{2/3}b^{2/3}$ and the integrated inverse scales by
$\ell/(\epsilon^{2/3}b^{2/3})=\epsilon^{-1/2}b^{-1}$.

\begin{proposition}[Exact pair tails]\label{prop:pair-tails}
As $\mu\to+\infty$ and $a\to+\infty$, respectively,
\begin{equation}\label{eq:pair-tails}
 \mathcal C_{\rm pair}(\mu)\sim\frac{C_0}{\sqrt2}\mu^{-3/4},
 \qquad
 \mathcal C_{\rm pair}(-a)\sim\frac\pi2 a^{-3/2}.
\end{equation}
Consequently
\begin{equation}\label{eq:pair-envelope}
 \mathcal C_{\rm pair}(\mu)\le C
 \begin{cases}(1+\mu)^{-3/4},&\mu\ge0,\\
 (1+|\mu|)^{-3/2},&\mu<0,
 \end{cases}
\end{equation}
and $\int_\R\mathcal C_{\rm pair}(\mu)^q\dd\mu<\infty$
exactly when $q>4/3$.
\end{proposition}
\begin{proof}
For positive $\mu$, put $r=\sqrt\mu$ and take neighborhoods
$|x\mp r|<\eta r$, where $\eta=\mu^{-3/8}$. On each one,
$(x^2-\mu)^2=4\mu(x\mp r)^2[1+O(\eta)]$, uniformly.
Their harmonic interval radii after scaling are comparable to
$\eta\mu^{3/4}\to\infty$. Bracketing and
\cref{lem:interval-response,prop:harmonic} therefore give the combined
contribution $2C_0(4\mu)^{-3/4}[1+o(1)]$.
For the upper bracket the omitted pieces contribute at most
\[
 \int_{|x-r|,|x+r|\ge\eta r}\frac{\dd x}{(x^2-\mu)^2}
 \le C\mu^{-3/2}\eta^{-1}=o(\mu^{-3/4}).
\]
To check this bound, set $x=\sqrt\mu\,z$; the reciprocal has integrable
$z^{-4}$ tails and only inverse-square singularities at $\pm1$,
cut off at distance $\eta$. Thus the positive tail is exact.
For $\mu=-a$, set $x=\sqrt a\,z$. The response becomes $a^{-3/2}$
times that for $-a^{-3}\partial_z^2+(1+z^2)^2$.
Dropping the derivative term bounds the latter by
$\int(1+z^2)^{-2}\dd z=\pi/2$. Conversely, the trial
$v=(1+z^2)^{-2}$ has finite derivative energy and gives
$\pi/2-a^{-3}\int|v'|^2$; it is admissible by energy cutoff.
This proves the second tail. Continuity and positivity on compact
sets now prove the envelope and the integrability assertion.
\end{proof}

The compact-parameter interval limit and these whole-line tails are
distinct statements. In particular, a rootless interval with
$\mu\to-\infty$ need not have the whole-line coefficient $\pi/2$ if
its radius divided by $\sqrt{|\mu|}$ stays bounded. No such joint
interval equivalent is used here.

\subsection{Uniform estimates for the compact cosine family}
Let $t=1-|c|$, $d=\epsilon^{1/3}$, and let
$\lambda_\epsilon(c)$ be the lowest eigenvalue of
$-\epsilon\partial_s^2+k_c$ on the circle.
\begin{lemma}[Uniform cosine localization]\label{lem:cosine-uniform}
There are constants independent of $|c|\le2$ and $0<\epsilon\le\epsilon_0$
such that
\begin{equation}\label{eq:cosine-bounds}
 \begin{array}{lll}
 t\ge0:&J_\epsilon(c)\le C\epsilon^{-1/4}(t+d)^{-3/4},&
 \lambda_\epsilon(c)\ge c_*\epsilon^{1/2}(t+d)^{1/2},\\[2pt]
 t<0:&J_\epsilon(c)\le C(|t|+d)^{-3/2},&
 \lambda_\epsilon(c)\ge c_*(|t|+d)^2.
 \end{array}
\end{equation}
With $b=1/2$, uniformly for $\mu=t/(\epsilon/2)^{1/3}$ in every
fixed compact set,
\begin{equation}\label{eq:cosine-fold-limit}
 J_\epsilon(c)=2\epsilon^{-1/2}
                    \{\mathcal C_{\rm pair}(\mu)+o(1)\}.
\end{equation}
On the zero-containing side there is also the following separate
uniform matching statement: along every sequence $|c|<1$ for which
$\epsilon/(1-|c|)^3\to0$,
\begin{equation}\label{eq:cosine-separated}
 J_\epsilon(c)=2C_0\epsilon^{-1/4}(1-c^2)^{-3/4}[1+o(1)].
\end{equation}
For fixed $|c|>1$ the response converges to $\int k_c^{-1}$.
\end{lemma}
\begin{proof}
It suffices to treat the fold $c=1$; translation by $\pi$ treats
$c=-1$. Center $x=s-\pi$ and choose a fixed small interval $|x|<x_0$.
The coordinate $y=2\sin(x/2)$ gives the exact identity
\[
 c+\cos s=y^2/2-t,\qquad
 \dd s=b_0(y)\dd y,\quad b_0(y)=(1-y^2/4)^{-1/2}.
\]
Here $b_0$ and $b_0^{-1}$ are bounded and positive. The derivative
energy has weight $b_0^{-1}$, and the source and reaction have weight
$b_0$. Scale $y=(4\epsilon)^{1/6}z$. For $\mu$ in a compact set the
interval expands, its weights converge locally uniformly to one,
and \cref{lem:interval-response} applies. The complementary fixed
wall interval has rate bounded away from zero and contributes $O(1)$.
This proves \eqref{eq:cosine-fold-limit}, the response bound
$C\epsilon^{-1/2}$, and, by the anchored inequality and Neumann
bracketing, the gap bound $c_*\epsilon^{2/3}$ when $|t|\le C_1d$
for any fixed $C_1$.

For $t\ge d$ near the fold, the roots in the $x$ coordinate are
$\pm r$, where $r\asymp\sqrt t$. On neighborhoods
$|x\mp r|<\eta_0r$ with a sufficiently small fixed $\eta_0>0$,
$k_c\asymp t(x\mp r)^2$; the two comparison constants are uniform.
The harmonic scale is $\ell\asymp(\epsilon/t)^{1/4}$ and
$r/\ell\asymp(t^3/\epsilon)^{1/4}$. When this ratio exceeds a fixed
large constant, \cref{lem:interval-response} bounds each Neumann
response by $C\epsilon^{-1/4}t^{-3/4}$ and its Rayleigh quotient below
by $c_*\sqrt{\epsilon t}$. On the remaining pieces,
$k_c\ge c_*t^2$ and $\int k_c^{-1}\le Ct^{-3/2}$;
the latter is at most $C\epsilon^{-1/4}t^{-3/4}$.
Also $t^2\ge\sqrt{\epsilon t}$, as $t^3\ge\epsilon$.
The bounded intermediate ratios are already covered by the compact
fold argument. This proves the inside bounds near the fold.
Away from the fold the roots have uniformly positive slopes and
uniformly separated neighborhoods, so the identical harmonic brackets
with fixed neighborhoods prove the inside bounds throughout $0\le t\le1$.

For $t=-\nu<0$, translation gives
$k_c=(\nu+1-\cos s)^2\ge\nu^2+(1-\cos s)^2$.
The quartic gap at $\nu=0$ thus gives
$\lambda_\epsilon(c)\ge\nu^2+c_*\epsilon^{2/3}$.
When $\nu\ge d$, the direct reciprocal-rate upper bound is
$\int k_c^{-1}\le C\nu^{-3/2}$, by comparison of $1-\cos s$
with squared distance to its minimum. For $\nu\le d$ use the compact
fold bound. This proves all of \eqref{eq:cosine-bounds}.

For sharp separated-root matching near a fold put
$\rho=\epsilon/t^3\to0$, and now choose shrinking relative
neighborhoods of radius $\eta r$, with $\eta=\rho^{1/8}$.
Taylor's formula at either root gives
$k_c=(1-c^2)(x\mp r)^2[1+O(\eta)]$ uniformly there.
Their scaled harmonic radii are comparable to
$\eta\rho^{-1/4}=\rho^{-1/8}\to\infty$.
The complementary reciprocal-rate integral is at most
$Ct^{-3/2}\eta^{-1}$, whose ratio to the proposed leading response
is $O(\rho^{1/4}/\eta)=O(\rho^{1/8})$.
Harmonic bracketing proves \eqref{eq:cosine-separated}.
Sequences remaining away from the folds are handled by fixed-root
localization, uniformly by compactness of their root neighborhoods.
These two alternatives prove the stated sequential uniformity.
Finally, for fixed $|c|>1$ the rate is positive and the reciprocal-rate
trial used in \cref{prop:harmonic} proves convergence.
\end{proof}

\subsection{Positive moments without design}
\begin{theorem}[Uniform-mobility disorder moments]\label{thm:uniform-moments}
For every fixed $q>0$, as $\epsilon\downarrow0$,
\begin{equation}\label{eq:uniform-moments}
 \E[J_\epsilon(c)^q]\sim
 \begin{cases}
 \displaystyle\frac{2^qC_0^q}{4}
 \Bfun\left(\frac12,1-\frac{3q}{4}\right)\epsilon^{-q/4},
                         &0<q<4/3,\\[6pt]
 \displaystyle\frac{2^{1/3}C_0^{4/3}}6
       \epsilon^{-1/3}\log(1/\epsilon),&q=4/3,\\[6pt]
 \displaystyle2^{q-4/3}\epsilon^{1/3-q/2}
       \int_\R\mathcal C_{\rm pair}(\mu)^q\dd\mu,&q>4/3.
 \end{cases}
\end{equation}
In particular,
\begin{align}
 \E J_\epsilon&\sim\frac{C_0^2}{\sqrt\pi}\epsilon^{-1/4},
                                  \label{eq:uniform-mean}\\
 \Var(J_\epsilon)&\sim2^{2/3}\epsilon^{-2/3}
                     \int_\R\mathcal C_{\rm pair}(\mu)^2\dd\mu.
                                  \label{eq:uniform-variance}
\end{align}
The mean coefficient is $12.18594958\ldots$. The variance coefficient
in \eqref{eq:uniform-variance} is defined by its convergent integral;
any decimal approximation to that integral is numerical evidence.
To express every formula at budget $M$, substitute $\epsilon=M/(2\pi)$,
including in the logarithm.
\end{theorem}
\begin{proof}
For $|c|<1$, \eqref{eq:cosine-separated} gives the pointwise limit
$\epsilon^{1/4}J_\epsilon(c)\to2C_0(1-c^2)^{-3/4}$;
for $|c|>1$ the limit is zero. On the inside,
\eqref{eq:cosine-bounds} bounds the $q$th power by $Ct^{-3q/4}$.
On the outside, writing $\nu=|t|=dz$, the same estimates give
$\epsilon^{1/4}J_\epsilon(c)\le C\nu^{-3/4}$:
indeed $z^{3/4}(1+z)^{-3/2}$ is bounded.
This integrable dominating function for $q<4/3$ proves the first
line of \eqref{eq:uniform-moments}; the density is $1/4$ and
$\int_{-1}^1(1-c^2)^{-3q/4}\dd c=
\Bfun(1/2,1-3q/4)$.

For $q>4/3$, use $t=(\epsilon/2)^{1/3}\mu$ near each of the two folds.
Equation \eqref{eq:cosine-fold-limit} and the rescaled envelope
\eqref{eq:cosine-bounds} permit dominated convergence on the expanding
$\mu$ intervals, by \cref{prop:pair-tails}. Contributions outside
fixed fold neighborhoods are $O(\epsilon^{-q/4}+1)$, smaller than
$\epsilon^{1/3-q/2}$. The two folds, density $1/4$, amplitude
$2^q\epsilon^{-q/2}$, and Jacobian $(\epsilon/2)^{1/3}$ give
$2^{q-4/3}$ exactly.

At $q=4/3$, \eqref{eq:cosine-separated} and $1-c^2=t(2-t)$ give
the inside separated-root contribution near either fold as
\[
 \frac14(2C_0)^{4/3}\epsilon^{-1/3}[t(2-t)]^{-1}
     [1+o(1)]\dd t.
\]
Here is a precise logarithmic cutoff argument. Fix $0<b<1/3$ and
a small fixed $t_0>0$. On $\epsilon^b\le t\le t_0$ the separated-root
error tends uniformly to zero. Its integral at both folds, divided by
$\epsilon^{-1/3}\log(1/\epsilon)$, tends to
$(2C_0)^{4/3}b/4$, with the factor $(2-t)^{-1}$ contributing its
limit $1/2$ to the logarithmic coefficient. The inside omitted interval
$0\le t\le\epsilon^b$ is bounded, using \eqref{eq:cosine-bounds}, by
\[
 C\epsilon^{-1/3}\int_0^{\epsilon^b}(t+d)^{-1}\dd t
 \le C\epsilon^{-1/3}
       \{1+(1/3-b)\log(1/\epsilon)\}.
\]
The outside-fold integral is $O(\epsilon^{-1/3})$, since
$\int_0^{t_0}(\nu+d)^{-2}\dd\nu\le d^{-1}$.
Regions away from folds have the same lower order. First let
$\epsilon\downarrow0$ and then $b\uparrow1/3$. The lower bracket
and the upper error above coincide in the limit, giving
$(2C_0)^{4/3}/12=2^{1/3}C_0^{4/3}/6$.
The beta--gamma identities give \eqref{eq:uniform-mean}.
For the variance, $(\E J_\epsilon)^2=O(\epsilon^{-1/2})$ is smaller
than the second moment, which has order $\epsilon^{-2/3}$.
\end{proof}

The squared coefficient of variation therefore grows as
$\epsilon^{-1/6}$ times a finite positive constant. This describes
variation between kinetic realizations, not non-Gaussian displacement
of a tracer. More explicitly, the almost-sure limiting random variable
$W=\lim_{\epsilon\downarrow0}\epsilon^{1/4}J_\epsilon(c)$ has an atom
of mass $1/2$ at zero, and its positive values are at least $2C_0$.
For $w\ge2C_0$, direct inversion of its expression above gives
\begin{equation}\label{eq:cosine-limit-law}
 \mathbb P(W>w)=\frac12\left[1-
           \sqrt{1-(2C_0/w)^{4/3}}\right]
 \sim2^{-2/3}C_0^{4/3}w^{-4/3}.
\end{equation}
The nonintegrability of higher moments of $W$ does not imply infinite
moments at fixed positive $\epsilon$; all those moments are finite.

\subsection{A uniform finite-bulk correction for this family}
The logarithmic remainder theorem already transfers divergent moments
of the uniform trial. The cosine family in fact has a stronger
realization-uniform expansion.
\begin{proposition}[Exchange flux and regular bulk limit]
\label{prop:uniform-bulk-limit}
Under the model assumptions of \cref{sec:model}, let
$h=H_{k_c,\epsilon}^{-1}1$ and $w=k_ch$. Uniformly for $|c|\le2$,
\begin{equation}\label{eq:uniform-flux}
 \|w-1\|_{L^2(\Gamma)}\le C\epsilon^{1/12}.
\end{equation}
Let $f_0$ be the mean-zero solution of
\begin{equation}\label{eq:regular-bulk-problem}
 -D_b\Delta f_0=u-V\quad\hbox{in }\Omega,\qquad
 D_b\partial_n f_0=-KV\quad\hbox{on }\Gamma,
\end{equation}
and set $\mathcal R_0=(D_b/Z)\int_\Omega|\nabla f_0|^2$.
Then
\begin{equation}\label{eq:uniform-bulk-expansion}
 \sup_{|c|\le2}\big|D_\flow(\epsilon,c)-\chi J_\epsilon(c)
                                     -\mathcal R_0\big|
 \le C\epsilon^{1/12}.
\end{equation}
Consequently every line of \eqref{eq:uniform-moments} transfers to
$\E[D_\flow^q]$ after multiplication by $\chi^q$; the mean and
variance equivalents transfer with factors $\chi$ and $\chi^2$.
\end{proposition}
\begin{proof}
The smooth periodic equation $-\epsilon h''+k_ch=1$ gives
$\int w=P$, the energy identity
$\epsilon\int|h'|^2+\int k_ch^2=J_\epsilon$, and
$\|h\|_2^2\le J_\epsilon/\lambda_\epsilon$.
Multiply the equation by $k_ch$ and integrate by parts. After using
$\int w=P$, the result is
\begin{equation}\label{eq:flux-identity}
 \|w-1\|_2^2+\epsilon\int k_c|h'|^2
                  =\frac\epsilon2\int k_c''h^2.
\end{equation}
With $g=c+\cos s$, direct differentiation gives
$k_c''=2(1-c^2)+6cg-4g^2$.
Discard its negative quadratic term and use
$\int|g|h^2\le(\int k_ch^2\int h^2)^{1/2}$ to obtain
\[
 \|w-1\|_2^2\le C\epsilon J_\epsilon
    \left\{\frac{(1-c^2)_+}{\lambda_\epsilon}
                                +\lambda_\epsilon^{-1/2}\right\}.
\]
For $t\ge0$, substitution from \eqref{eq:cosine-bounds} bounds this by
$C[\epsilon^{1/4}t(t+d)^{-5/4}+\epsilon^{1/2}(t+d)^{-1}]
\le C\epsilon^{1/6}$. For $t<0$ the first term vanishes and the
bound is $C\epsilon(|t|+d)^{-5/2}\le C\epsilon^{1/6}$.
This proves \eqref{eq:uniform-flux}.

The compatibility of \eqref{eq:regular-bulk-problem} is
$\int_\Omega(u-V)=KPV$. Bulk Poincar\'e and the trace estimate,
together with \eqref{eq:uniform-flux}, show that the load
$\mathcal L_\epsilon$ in \eqref{eq:bulk-remainder} differs in the
energy-dual norm from
$\mathcal L_0(f)=\int_\Omega(u-V)f-KV\int_\Gamma f_\Gamma$
by $O(\epsilon^{1/12})$, uniformly in $c$.
Dropping the nonnegative Schur penalty gives the upper bound
$\mathcal R\le\mathcal R_0+C\epsilon^{1/12}$: the unpenalized
supremum is the squared dual norm of its load divided by $Z$.
On the stated $C^2$ domain the Neumann solution $f_0$ is in $H^2$
for $u\in L^2$ and constant boundary datum, by the standard Neumann
regularity theorem \citep[Section~2.4]{Grisvard2011}. Its trace has a square
integrable tangential derivative. In the Schur infimum choose
$v=(f_0)_\Gamma$, obtaining
$\mathcal S_\epsilon((f_0)_\Gamma)
 \le\epsilon\int_\Gamma|\partial_s(f_0)_\Gamma|^2$.
Inserting $f_0$ in the bulk supremum gives the matching lower bound
$\mathcal R\ge\mathcal R_0-C(\epsilon^{1/12}+\epsilon)$.
This proves \eqref{eq:uniform-bulk-expansion}.

The correction is uniformly bounded while every scalar positive moment
diverges. For $0<q\le1$ use $|(a+b)^q-a^q|\le b^q$ for $a,b\ge0$.
For $q>1$ the difference is bounded by $C_q(J_\epsilon^{q-1}+1)$;
H\"older bounds its expectation by a lower-order quantity than
$\E J_\epsilon^q$. For variance use the second-moment equivalent and
$\E D_\flow=O(\epsilon^{-1/4})$. The same statements permit an
additional uniformly bounded nonnegative molecular diffusivity.
\end{proof}

\subsection{Why local zero statistics and mean rates do not suffice}
Two counterexamples delimit the averaging operation. First, on the same
circle take
\begin{equation}\label{eq:gaussian-collapse}
 g(s)=\xi_1\cos s+\xi_2\sin s,\qquad
 \xi_1,\xi_2\ \hbox{independent }N(0,1),\qquad k=g^2.
\end{equation}
At each point $(g,g')$ is a nondegenerate standard Gaussian vector.
Writing $R=(\xi_1^2+\xi_2^2)^{1/2}$, there are almost surely two
simple zeros with slope magnitude $R$. The local limiting zero sum
is $2R^{-3/2}$, with finite mean because $R$ has density
$r e^{-r^2/2}$. Its second moment is infinite. Nevertheless, the
constant trial in \eqref{eq:J-quotient} gives, for \emph{every}
nonnegative integrable mobility field, even chosen after observing $g$,
\begin{equation}\label{eq:gaussian-response-lower}
 J_{g^2}(D)\ge\frac{P^2}{\int g^2}=\frac{2P}{R^2},
 \qquad \E J_{g^2}(D)=\infty.
\end{equation}
Thus a finite mean samplewise small-$\epsilon$ limit does not justify
interchanging that limit with expectation. The missing control here
is on nearly vanishing amplitudes of the entire rate field. This
family does not satisfy the uniform anchoring assumption
\eqref{eq:rate-assumptions}.

For context, the classical first marked Kac--Rice formula weights a
zero's slope by its magnitude \citep[Theorems~6.2 and~6.4]{AzaisWschebor2009}.
More precisely, suppose a
stationary Gaussian process has almost surely $C^1$ paths, locally
finite simple zeros, mean $m_g$, variances $\sigma_0^2>0$ and
$\sigma_1^2>0$ for its value and derivative, on a neighborhood of the
observation interval. These hypotheses permit
the weighted-root formula with the derivative as the auxiliary Gaussian
field. Bounded continuous truncations followed by monotone convergence
give the identity for the negative-power marks used below:
\begin{equation}\label{eq:marked-rice}
 \E\sum_{g(s_j)=0}\psi(|g'(s_j)|)
  =L p_g(0)\E[|U|\psi(|U|)],\qquad U\sim N(0,\sigma_1^2),
\end{equation}
where the sum is on an interval of length $L$ and
$p_g(0)=e^{-m_g^2/(2\sigma_0^2)}/(\sqrt{2\pi}\sigma_0)$.
Stationarity makes $g$ and $g'$ independent at one point.
For the zero sum $S_L=\sum|g'(s_j)|^{-3/2}$, on any interval with $L>0$, this yields
\[
 \E S_L=L p_g(0)\sigma_1^{-1/2}
                 \frac{2^{-1/4}\Gamma(1/4)}{\sqrt\pi}<\infty,
 \qquad
 \E S_L^2\ge Lp_g(0)\E|U|^{-2}=\infty.
\]
Equivalently, the slope magnitude at a typical zero, sampled by zero
intensity, has density $f_{\rm zero}$, with
\[
 f_{\rm zero}(r)=r\sigma_1^{-2}e^{-r^2/(2\sigma_1^2)},\qquad
 \mathbb P_{\rm zero}(R^{-3/2}>w)
       =1-\exp[-w^{-4/3}/(2\sigma_1^2)].
\]
The inverse-$3/2$ mark has finite $q$th moment exactly for $q<4/3$. These are applications of the classical identity, not
new Gaussian-process results. They neither control global amplitude
collapse nor provide independence of nearby zeros or a stable limit.

Second, in the bounded cosine ensemble,
$\E k_c(s)=4/3+\cos^2s\ge4/3$. Replacing the realized rate by this
mean gives
\begin{equation}\label{eq:mean-rate-failure}
 J_{\E k_c}(D)\le\int_\Gamma\frac{\dd s}{\E k_c(s)}\le\frac{3P}{4}
\end{equation}
for every $D\ge0$, by discarding derivative energy. The actual
uniform-design mean instead diverges as \eqref{eq:uniform-mean}.
Averaging the rate before solving the source problem therefore removes
the kinetic defects responsible for the singular response.
